\begin{lemma}
Let $\mathcal H$ be a $3$-uniform $2$-simple hypergraph of maximum degree at
most $D$, and let $e\in E(\mathcal H)$ intersect exactly $3(D-1)$ other
edges.  Put
\[
X(e,\mathcal H)=
\left(\bigcup\{h\in E(\mathcal H):h\cap e\ne\emptyset\}\right)\setminus e
\]
and, for $v\in e$ and $x\in X(e,\mathcal H)$, let
$\deg_{\mathcal H}(v,x)$ be the number of edges containing both $v$ and $x$.
Then:
\begin{enumerate}
\item every edge $h\ne e$ satisfies $|h\cap e|\le 1$;
\item for every $x\in X(e,\mathcal H)$,
$\sum_{v\in e}\deg_{\mathcal H}(v,x)\le D$;
\item for every $v\in e$,
$\sum_{x\in X(e,\mathcal H)}\deg_{\mathcal H}(v,x)=2(D-1)$.
\end{enumerate}
\end{lemma}
\begin{proof}
The hypotheses are the $3$-uniform and maximum-degree components of the
local $H(3,2,D)$ setup described by \noderef{HypergraphClass}.  By
\noderef{UniformHypergraph}, the edge $e$ has exactly three vertices; write
$e=\{v_1,v_2,v_3\}$.

For the first assertion, suppose some edge $h\ne e$ contains two
vertices of $e$, say $v_1$ and $v_2$.  Counting neighboring edges of $e$ by
the vertex of $e$ that they contain gives at most
$3(D-1)$ incidences: by \noderef{MaxDegreeAtMost} and the incidence-counting
definition of degree in \noderef{HypergraphDegree}, each $v_i$ belongs to at
most $D$ edges of $\mathcal H$, hence to at most $D-1$ edges other than
$e$.  The edge $h$ is counted at least twice, so the number of distinct edges
meeting $e$ is at most $3(D-1)-1$, contradicting the hypothesis that there
are exactly $3(D-1)$ such edges.  Hence every edge other than $e$ contains at
most one vertex of $e$.

Fix $x\in X(e,\mathcal H)$.  By the first assertion, the sets of edges
counted by $\deg_{\mathcal H}(v_i,x)$ for $i=1,2,3$ are pairwise disjoint:
an edge counted for two different $v_i$ would contain $x$ and two vertices
of $e$.  Their union is a subset of the set of all edges containing $x$.
Using again \noderef{HypergraphDegree} to interpret $\deg_{\mathcal H}(x)$
as the number of incident edges and \noderef{MaxDegreeAtMost} to bound that
degree, we get
\[
\sum_{i=1}^3\deg_{\mathcal H}(v_i,x)\le \deg_{\mathcal H}(x)\le D.
\]

Finally fix $v_i\in e$.  The same incidence count used above has total
$3(D-1)$ and is a sum of three quantities, each at most $D-1$.  Since the
number of distinct neighboring edges is also $3(D-1)$ and no edge is counted
twice, each $v_i$ is contained in exactly $D-1$ neighboring edges.  Each such
edge has exactly three vertices by \noderef{UniformHypergraph} and meets
$e$ only in $v_i$, so it contains exactly two
vertices from $X(e,\mathcal H)$ and contributes exactly two to
$\sum_{x\in X(e,\mathcal H)}\deg_{\mathcal H}(v_i,x)$.  Thus the sum is
$2(D-1)$.
\end{proof}
